%!TEX root = TheoretischePhysikA.tex

\section{Vorlesung: Starrer Körper}
Ein Vielteilchensystem mit festen Abständen zwischen Teilchen wird als \emph{starrer Körper} bezeichnet.
\figbox
$3N$ Zahlen werden benötigt um ein allgemeines Vielteilchensystem zu beschreiben. Wegen der Zwangsbedingungen, dass
\begin{equation}
  \|\mbf{r}_{jk}\| = \|\mbf{r}_{j}-\mbf{r}_{k}\| = \text{Konst.,} \quad \forall j\not= k,
\end{equation}
sind allerdings nur $6$ unabhängige Koordinaten übrig. Ein starrer Körper ist durch Angabe von den folgenden Größen beschrieben:
\begin{enumerate}
  \item $3$ Zahlen für die Schwerpunktsposition $\mbf{R}(t)$ und
  \item $3$ Zahlen für die Drehung relativ zum Bezugssystem (z.B.\ Eulerwinkel.)
\end{enumerate}
Für den Gesamtimpuls und den Drehimpuls gilt:
\begin{equation}
  \eqbox{\left(\frac{d\mbf{P}}{dt}\right)_{\Sigma} =\mbf{F}^{\text{ext.}}\quad\text{und}\quad \left(\frac{d\mbf{L}}{dt}\right)_{\Sigma} =\mbf{M}^{\text{ext.}}.}
\end{equation}

Wir führen ein zweites Bezugssystem $\Sigma'$ ein, das sich mit dem Körper rotiert. Das heisst, $(\mbf{v}_j)_{\Sigma'} = \mbf{0}$. Der Gesamtdrehimpuls in $\Sigma$ lautet
\begin{equation}
  \begin{split}
    \mbf{L} &= \sum_{j=1}^N m_j\mbf{r}_j\times (\mbf{v}_j)_{\Sigma} \\
    &= \sum_{j=1}^N m_j\mbf{r}_j\times [(\mbf{v}_j)_{\Sigma'} + \boldsymbol{\omega}\times \mbf{r}_j)] \\
    &= \sum_{j=1}^N m_j\mbf{r}_j\times (\boldsymbol{\omega}\times \mbf{r}_j) \\
    &= \sum_{j=1}^N m_j[(\mbf{r}_j\cdot \mbf{r}_j)\boldsymbol{\omega} -  (\mbf{r}_j\cdot \boldsymbol{\omega}) \mbf{r}_j].
  \end{split}
\end{equation}
Man kann dieses Gleichungssystem als eine Matrixgleichung darstellen:
\begin{equation}
  \begin{pmatrix}
    L_x \\ L_y \\ L_z
  \end{pmatrix} = \sum_{j=1}^N m_j
  \begin{pmatrix}
    y_j^2 + z_j^2 & -x_jy_j & -x_jz_j \\
    -x_jy_j & x_j^2 + z_j^2 & -y_jz_j \\
    -x_jz_j & -y_jz_j &  x_j^2 + y_j^2
  \end{pmatrix}
  \begin{pmatrix}
    \omega_x \\ \omega_y \\ \omega_z
  \end{pmatrix}.
\end{equation}
Wir definieren dazu die Matrix
\begin{equation}
  \mbf{J}\equiv\sum_{j=1}^N m_j
  \begin{pmatrix}
    y_j^2 + z_j^2 & -x_jy_j & -x_jz_j \\
    -x_jy_j & x_j^2 + z_j^2 & -y_jz_j \\
    -x_jz_j & -y_jz_j &  x_j^2 + y_j^2
  \end{pmatrix}
\end{equation}
und erhalten die Matrixgleichung
\begin{equation}
  \eqbox{\mbf{L} = \mbf{J}\boldsymbol{\omega}.}
\end{equation}

Die symmetrische Matrix $\mbf{J}$ hat die folgenden Eigenschaften:
\begin{enumerate}
  \item $\mbf{J}$ hat nicht-negative Eigenwerte $(J_x,J_z,J_y)$.
  \item Es gibt ein Drehmatrix $\mbf{D}$, so dass
  \begin{equation}
    \mbf{J} = \mbf{D}\boldsymbol{\Delta}\mbf{D}^{-1} \quad \text{mit}\quad \boldsymbol{\Delta}=\begin{pmatrix}
      J_x & 0 & 0 \\
      0 & J_y & 0 \\
      0 & 0 & J_z
    \end{pmatrix}.
  \end{equation}
  \item Die Drehmatrix gibt die Richtungen der Hauptträgheitsachsen an:
  \begin{equation}
    \mbf{j}_x \equiv \mbf{D}\hat{\mbf{x}}, \quad \mbf{j}_y \equiv \mbf{D}\hat{\mbf{y}} \quad \text{und}\quad \mbf{j}_z \equiv \mbf{D}\hat{\mbf{z}}.
  \end{equation}
  Für Rotationen um die Haupträgheitsachsen gilt
  \begin{equation}
    \mbf{L} \| \boldsymbol{\omega}.
  \end{equation}  
\end{enumerate}

Für eine beliebige Rotation mit Winkelgeschwindigkeitsvektoren
\begin{equation}
  \boldsymbol{\omega} = \omega_x \mbf{j}_x +\omega_y \mbf{j}_y +\omega_z \mbf{j}_z
\end{equation}
ist 
\begin{equation}
  \begin{split}
    \mbf{L} &= \mbf{J}\boldsymbol{\omega}\\
    &= \omega_x \mbf{J}\mbf{j}_x +\omega_y \mbf{J}\mbf{j}_y +\omega_z \mbf{J}\mbf{j}_z \\
    &= \omega_x J_x\mbf{j}_x +\omega_y J_y\mbf{j}_y +\omega_z J_z\mbf{j}_z.
  \end{split}
\end{equation}
Im Körperfesten rotierenden Bezugssystem gilt
\begin{equation}
  \left(\frac{d\mbf{L}}{dt}\right)_{\Sigma}  + \boldsymbol{\omega}\times \mbf{L} = \mbf{M}.
\end{equation}
Im Hauptachsensystem lautet diese Gleichung
\begin{equation}
  \begin{pmatrix}
    \dot{\omega}_x J_x \\
    \dot{\omega}_y J_y \\
    \dot{\omega}_z J_z 
  \end{pmatrix}
  +
  \begin{pmatrix}
    \omega_x \\
    \omega_y \\
    \omega_z 
  \end{pmatrix}
  \times 
  \begin{pmatrix}
    \omega_x J_x \\
    \omega_y J_y \\
    \omega_z J_z
  \end{pmatrix}
  =
  \begin{pmatrix}
    M_x \\
    M_y \\ 
    M_z
  \end{pmatrix}.
\end{equation}
In Komponentenform wird diese Matrixgleichung durch 
\begin{equation}
  \begin{split}
    \dot{\omega}_x J_x + \omega_y\omega_z(J_z-J_y) &= M_x \\
    \dot{\omega}_y J_y + \omega_x\omega_z(J_x-J_z) &= M_y \\
    \dot{\omega}_z J_z + \omega_x\omega_y(J_y-J_x) &= M_z \\
  \end{split}
\end{equation}
gegeben. Dieses System wird als die \emph{Euler'sche Gleichungen} bezeichnet und stellt die Bewegungsgleichung für die Drehbewegung eines starren Körpers dar. 

\begin{beispiel}
  Angenommen es sei ein Körper mit $J_j\not=J_k$, $j\not=k$, zum Beispiel ein Quader. Der Körper bewege sich mit konstanten Winkelgeschwindigkeit $\boldsymbol{\omega}=\mbf{0}$ ohne äußere Drehmoment ($\mbf{M}^{\text{ext}}=\mbf{0}$). Die Euler'schen Gleichungen lauten 
  \begin{equation}
    \begin{split}
    \omega_y\omega_z(J_z-J_y) &= 0 \\
    \omega_x\omega_z(J_x-J_z) &= 0 \\
    \omega_x\omega_y(J_y-J_x) &= 0. \\
  \end{split}
  \end{equation}
  Es gibt drei Lösungen:
  \begin{enumerate}
    \item $\omega_x = \text{Konst.}$ und $\omega_y=\omega_z  = 0$.
    \item $\omega_y = \text{Konst.}$ und $\omega_x=\omega_z  = 0$.
    \item $\omega_z = \text{Konst.}$ und $\omega_x=\omega_y  = 0$.
  \end{enumerate}
  Dies entspricht Rotationen um die Hauptträgheitsachsen.
\end{beispiel}

\begin{beispiel}(Symmetrischer Kreisel)
  In diesem Fall wird $J_x = J_y \not= J_z$. Die Euler'schen Gleichungen lauten
  \begin{equation}
    \begin{split}
    \dot{\omega}_x J_x + \omega_y\omega_z(J_z-J_y) &= 0 \\
    \dot{\omega}_y J_y + \omega_x\omega_z(J_x-J_z) &= 0 \\
    \dot{\omega}_z J_z  &= 0. \\
  \end{split}
  \end{equation}
  Zum ersten stellen wir fest, dass $\omega_z = \text{Konst.}$. Die anderen Komponenten erfüllen die Differentialgleichungen
  \begin{equation}
    \dot{\omega}_x = \Omega \omega_y\quad\text{und}\quad \dot{\omega}_y = \Omega \omega_x,
  \end{equation}
  wobei $\Omega = \omega_z(J_x-J_z)/J_x$. Die Lösung lautet
  \begin{equation}
    \begin{pmatrix}
      \omega_x(t) \\
      \omega_y(t) \\
      \omega_z(t) 
    \end{pmatrix}
    =
    \begin{pmatrix}
      \omega_0 \sin(\Omega t) \\
      \omega_0 \cos(\Omega t) \\
      0
    \end{pmatrix}.
  \end{equation}
\end{beispiel}